\scanpage{56}
La démonstration de 1.5. est de la théorie d'intersections standart, que nous nous dispensons ici d\deleted{\unsure{e}}'expliciter.

Suivant du point de vue cohomologie $l$-adique la construction qu'on vient de faire, on vérifie sans difficulté :

\uline{Théorème} 1.6. Avec les notations du th. 1.1. considérons $\gamma^n(Z)\in H^{2n}(X\times Y)(n)$, et l'homomorphisme correspondant
\[
a(\gamma^n(Z)):H^1(X)\longrightarrow H^1(Y)
\]
\deleted{Considérons d'autre} Considérons d'autre part \deleted{Alors on a} les h\deleted{\unsure{iso|om}}om\inserted{\unsure{o|à}}morphismes canoniques (isom. mod \deleted{\unsure{torsion})}\inserted{groupes finis}) $\varphi_X^*:H^1(A_X)\unsure{-|\longrightarrow} H^1(X)$, $\varphi_Y^*:H^1(A_Y)-H^1(Y)$. On a commutativité dans le diagramme d'homomorphismes
\diagram{\begin{tikzcd}[ampersand replacement=\&]
H^1(X)\arrow[r,"a(\gamma^n(Z))"] \& H^1(Y)\\
H^1(A_X)\arrow[u,"\varphi_X^*"']\arrow[r,"u_Z^{(1)}"] \& H^1(A_Y)\arrow[u,"\varphi_Y^*"']
\end{tikzcd}}
(Les \deleted{\unsure{f}}cohomologies sont les cohomologies à coefficient $\mathrm{Z}_{\ell}$). \deleted{Ce diagramme est \uline{commutatif}}.
% archive sheet number: 55
